\documentclass[12pt,a4paper]{article}

\usepackage{geometry}
\geometry{a4paper, top=2.7cm, bottom=2.7cm, left=2.3cm, right=2.3cm}

\usepackage{graphicx}
\graphicspath{{./}{pics/}}
\usepackage{float}
\usepackage{caption}
\usepackage{subcaption}
\usepackage{booktabs}
\usepackage{array}
\usepackage{longtable}
\usepackage{enumitem}
\usepackage[table]{xcolor}
\usepackage{amsmath}
\usepackage[hidelinks]{hyperref}

\usepackage{xepersian}

\definecolor{headerbgA}{RGB}{47, 85, 151}
\definecolor{rowlightA}{RGB}{237, 242, 250}
\definecolor{headerbgB}{RGB}{40,60,90}
\definecolor{rowlightB}{RGB}{235,240,247}

\IfFontExistsTF{XB Niloofar}
  {\settextfont[Scale=1.0]{XB Niloofar}}
  {\settextfont[Scale=1.0]{Amiri}}

\IfFontExistsTF{Times New Roman}
  {\setlatintextfont{Times New Roman}}
  {\setlatintextfont{Latin Modern Roman}}

\IfFontExistsTF{XB Niloofar}
  {\setdigitfont{XB Niloofar}}
  {}

\linespread{1.35}
\setlength{\parskip}{4pt}

\newcommand{\en}[1]{\lr{#1}}

\captionsetup{font=footnotesize,labelfont=bf}

\title{}
\author{}
\date{}

\begin{document}

\begin{titlepage}
\begin{center}
\vspace*{2cm}

\vspace{2cm}

{\huge \textbf{گزارش آزمایشگاه مدارهای منطقی}}

\vspace{0.6cm}

{\LARGE \textbf{آزمایش پنجم: واحد محاسبات و منطق}\ \lr{(ALU)}}

\vspace{2.2cm}

{\Large علی‌اصغر طایری -- \lr{404171133}}

\vspace{0.3cm}

{\Large رضا حضرتی -- \lr{404105767}}

\vfill

\end{center}
\end{titlepage}

\tableofcontents
\newpage

\part{بخش اول: آشنایی با تراشه‌ی \lr{74181} و طراحی مدار ثبات‌ها و کنترل‌کننده}

\section{مقدمه و تعریف مسئله}

هدف از این آزمایش، آشنایی عملی با واحد محاسبات و منطق \lr{(Arithmetic Logic Unit - ALU)} و بررسی نحوه‌ی عملکرد تراشه‌ی معروف \lr{74181} است. این تراشه یکی از قدیمی‌ترین و شناخته‌شده‌ترین \lr{ALU}های تجاری است که قابلیت انجام عملیات محاسباتی (مانند جمع) و عملیات منطقی (مانند \lr{NOT}، \lr{AND} و غیره) را بر روی دو عملوند ۴ بیتی دارد و انتخاب نوع عملیات از طریق ورودی‌های انتخاب \lr{S3-S0} و ورودی حالت \lr{M} صورت می‌گیرد.

در این آزمایش قرار است مداری طراحی و در نرم‌افزار \lr{Proteus} پیاده‌سازی شود که دارای دو ثبات داده‌ی ۴ بیتی \lr{A} و \lr{B} و یک واحد \lr{ALU} است. مقدار این دو ثبات از طریق دو مالتی‌پلکسر قابل انتخاب است تا هم بتوان مقدار جدیدی از خطوط ورودی \lr{D0-D3} در آن‌ها بارگذاری کرد و هم بتوان خروجی \lr{ALU} را مجدداً به یکی از ثبات‌ها بازگرداند (\lr{feedback}). یک مدار کنترل‌کننده (\lr{Controller}) نیز بر اساس سه ورودی دستور \lr{M0-M2} مشخص می‌کند که در هر لحظه چه عملیاتی روی محتوای ثبات‌ها انجام شود.

جدول عملیات موردنظر که مبنای طراحی کنترل‌کننده قرار گرفته، در جدول \ref{tab:ops} آمده است.

\begin{table}[H]
\centering
\caption{عملیات صورت‌گرفته در مدار برحسب ورودی‌های \lr{M2-M0}}
\label{tab:ops}
\renewcommand{\arraystretch}{1.4}
\rowcolors{2}{rowlightA}{white}
\begin{tabular}{>{\centering\arraybackslash}p{1cm} >{\centering\arraybackslash}p{1cm} >{\centering\arraybackslash}p{1cm} | >{\centering\arraybackslash}p{4cm}}
\toprule
\rowcolor{headerbgA}
\textcolor{white}{\lr{M2}} & \textcolor{white}{\lr{M1}} & \textcolor{white}{\lr{M0}} & \textcolor{white}{عملیات} \\
\midrule
0 & 0 & 0 & \lr{$A \leftarrow D_3D_2D_1D_0$} \\
0 & 0 & 1 & \lr{$B \leftarrow D_3D_2D_1D_0$} \\
0 & 1 & 0 & \lr{$A \leftarrow A$} \\
0 & 1 & 1 & \lr{$A \leftarrow B$} \\
1 & 0 & 0 & \lr{$clear(A)$} \\
1 & 0 & 1 & \lr{$A \leftarrow not(A)$} \\
1 & 1 & 0 & \lr{$A \leftarrow and(A,B)$} \\
1 & 1 & 1 & \lr{$A \leftarrow add(A,B)$} \\
\bottomrule
\end{tabular}
\end{table}

مسئله‌ی اصلی این آزمایش، طراحی بلوک کنترل‌کننده‌ای است که با دریافت سه بیت \lr{M2-M0}، سیگنال‌های انتخاب مالتی‌پلکسرها (\lr{SA} و \lr{SB}) و سیگنال‌های انتخاب تابع \lr{ALU} (یعنی \lr{S3-S0} و ورودی \lr{M}) را به‌گونه‌ای تولید کند که هشت عملیات جدول بالا به‌درستی روی ثبات‌های \lr{A} و \lr{B} اعمال شوند.

\section{بررسی و انتخاب روش طراحی}

برای پیاده‌سازی این مدار، از سه تراشه‌ی اصلی زیر استفاده شد:

\begin{itemize}
\item \lr{IC 74181}: واحد محاسبات و منطق ۴ بیتی که عملیات محاسباتی و منطقی مورد نیاز جدول \ref{tab:ops} را انجام می‌دهد.
\item \lr{IC 74175}: چهار فلیپ‌فلاپ \lr{D} در یک بسته، که برای ساخت دو ثبات ۴ بیتی \lr{A} و \lr{B} از دو عدد از این تراشه استفاده شد.
\item \lr{IC 74157}: مالتی‌پلکسر چهارگانه‌ی ۲ به ۱، که برای انتخاب ورودی هر ثبات (بین داده‌ی خارجی و مقدار بازخوردی) به کار رفت.
\end{itemize}

با بررسی جدول تابع \lr{74181} (که در دیتاشیت این تراشه موجود است)، مشخص شد برای تولید هر یک از توابع موردنیاز، ترکیب مشخصی از ورودی حالت \lr{M} و ورودی‌های انتخاب \lr{S3S2S1S0} لازم است. جدول \ref{tab:alu-func} این نگاشت را نشان می‌دهد.

\begin{table}[H]
\centering
\caption{کدهای انتخاب \lr{ALU} برای توابع مورد نیاز (طبق دیتاشیت \lr{74181})}
\label{tab:alu-func}
\renewcommand{\arraystretch}{1.4}
\rowcolors{2}{rowlightA}{white}
\begin{tabular}{>{\centering\arraybackslash}p{3.5cm} >{\centering\arraybackslash}p{0.8cm} >{\centering\arraybackslash}p{1.8cm} >{\centering\arraybackslash}p{0.8cm} >{\centering\arraybackslash}p{2.8cm}}
\toprule
\rowcolor{headerbgA}
\textcolor{white}{تابع موردنیاز} & \textcolor{white}{\lr{M}} & \textcolor{white}{\lr{S3S2S1S0}} & \textcolor{white}{\lr{Cn}} & \textcolor{white}{توضیح} \\
\midrule
\lr{$F=A$} (نگه‌داری \lr{A}) & 1 & 1111 & - & عبور منطقی \lr{A} \\
\lr{$F=B$} (انتقال \lr{B} به \lr{A}) & 1 & 1010 & - & عبور منطقی \lr{B} \\
\lr{$F=0$} (\lr{clear}) & 1 & 0011 & - & صفر منطقی \\
\lr{$F=\overline{A}$} (\lr{not}) & 1 & 0000 & - & \lr{NOT A} \\
\lr{$F=A \cdot B$} (\lr{and}) & 1 & 1011 & - & \lr{AND} منطقی \\
\lr{$F=A+B$} (\lr{add}) & 0 & 1001 & 1 & جمع محاسباتی \\
\bottomrule
\end{tabular}
\end{table}

با توجه به این جدول، روش طراحی انتخاب‌شده به این صورت است: کنترل‌کننده یک مدار ترکیبی (کدکننده) است که ورودی سه‌بیتی \lr{M2-M0} را به خروجی‌های \lr{SA}، \lr{SB} و \lr{S3-S0} (به‌همراه ورودی \lr{M} و \lr{Cn} تراشه‌ی \lr{ALU}) ترجمه می‌کند. از آنجا که تعداد حالت‌های ورودی محدود (۸ حالت) و خروجی‌ها نیز مشخص هستند، ساده‌ترین و قابل‌اطمینان‌ترین روش، استخراج مستقیم جدول درستی کنترل‌کننده از روی جدول \ref{tab:ops} و \ref{tab:alu-func} و سپس پیاده‌سازی آن با گیت‌های منطقی پایه (\lr{AND}، \lr{OR}، \lr{NOT}) بود؛ این روش نسبت به استفاده از \lr{ROM} یا حافظه، برای مداری با این حجم محدود، از نظر تعداد قطعات و سادگی سیم‌کشی در \lr{Proteus} به‌صرفه‌تر تشخیص داده شد.

\section{بلوک‌دیاگرام طرح پیشنهادی اولیه}

بلوک‌دیاگرام کلی مدار پیشنهادی، مطابق شکل داده‌شده در راهنمای آزمایش، شامل دو مالتی‌پلکسر ۴ بیتی در ورودی دو ثبات \lr{A} و \lr{B}، خود دو ثبات \lr{A} و \lr{B}، واحد \lr{ALU} و بلوک کنترل‌کننده است.

\begin{itemize}
\item مالتی‌پلکسر \lr{A}: بین خطوط ورودی خارجی \lr{D3-D0} و خروجی \lr{F} تراشه‌ی \lr{ALU} (بازخورد) یکی را، بسته به \lr{SA}، انتخاب کرده و به ورودی ثبات \lr{A} می‌دهد.
\item مالتی‌پلکسر \lr{B}: بین خطوط ورودی خارجی \lr{D3-D0} و خروجی خودِ ثبات \lr{B} (برای حفظ مقدار فعلی در حالت‌هایی که \lr{B} نباید تغییر کند)، بسته به \lr{SB}، یکی را انتخاب کرده و به ورودی ثبات \lr{B} می‌دهد.
\item ثبات‌های \lr{A} و \lr{B}: با یک پالس ساعت مشترک (\lr{clock}) در هر لبه‌ی ساعت، مقدار ورودی مالتی‌پلکسر مربوط به خود را ذخیره می‌کنند.
\item واحد \lr{ALU}: مقادیر ذخیره‌شده در \lr{A} و \lr{B} را دریافت کرده و بسته به \lr{S3-S0} و \lr{M}، خروجی \lr{F} را تولید می‌کند.
\item کنترل‌کننده: با دریافت \lr{M2-M0}، سیگنال‌های \lr{SA}، \lr{SB} و \lr{S3-S0} را طبق جدول \ref{tab:controller} تولید می‌کند.
\end{itemize}

\begin{table}[H]
\centering
\caption{جدول درستی کنترل‌کننده}
\label{tab:controller}
\renewcommand{\arraystretch}{1.4}
\rowcolors{2}{rowlightA}{white}
\begin{tabular}{>{\centering\arraybackslash}p{0.7cm} >{\centering\arraybackslash}p{0.7cm} >{\centering\arraybackslash}p{0.7cm} | >{\centering\arraybackslash}p{0.7cm} >{\centering\arraybackslash}p{0.7cm} | >{\centering\arraybackslash}p{0.7cm} | >{\centering\arraybackslash}p{0.7cm} >{\centering\arraybackslash}p{0.7cm} >{\centering\arraybackslash}p{0.7cm} >{\centering\arraybackslash}p{0.7cm}}
\toprule
\rowcolor{headerbgA}
\textcolor{white}{\lr{M2}} & \textcolor{white}{\lr{M1}} & \textcolor{white}{\lr{M0}} & \textcolor{white}{\lr{SA}} & \textcolor{white}{\lr{SB}} & \textcolor{white}{\lr{M}} & \textcolor{white}{\lr{S3}} & \textcolor{white}{\lr{S2}} & \textcolor{white}{\lr{S1}} & \textcolor{white}{\lr{S0}} \\
\midrule
0 & 0 & 0 & 0 & 1 & \lr{X} & \lr{X} & \lr{X} & \lr{X} & \lr{X} \\
0 & 0 & 1 & \lr{X} & 0 & \lr{X} & \lr{X} & \lr{X} & \lr{X} & \lr{X} \\
0 & 1 & 0 & 1 & 1 & 1 & 1 & 1 & 1 & 1 \\
0 & 1 & 1 & 1 & 1 & 1 & 1 & 0 & 1 & 0 \\
1 & 0 & 0 & 1 & 1 & 1 & 0 & 0 & 1 & 1 \\
1 & 0 & 1 & 1 & 1 & 1 & 0 & 0 & 0 & 0 \\
1 & 1 & 0 & 1 & 1 & 1 & 1 & 0 & 1 & 1 \\
1 & 1 & 1 & 1 & 1 & 0 & 1 & 0 & 0 & 1 \\
\bottomrule
\end{tabular}
\end{table}

\noindent در سطرهای اول و دوم که ثبات مقصد فقط بارگذاری از ورودی خارجی است، سایر خروجی‌ها اهمیتی ندارند و در جدول با \lr{X} (بی‌اهمیت) نشان داده شده‌اند.

شکل \ref{fig:main} بلوک‌دیاگرام کلی و شماتیک نهایی پیاده‌سازی‌شده در \lr{Proteus} را نشان می‌دهد.

\begin{figure}[H]
\centering
\includegraphics[width=0.95\textwidth]{5-1-main.png}
\caption{شماتیک کلی مدار پیاده‌سازی‌شده در \lr{Proteus}}
\label{fig:main}
\end{figure}

\section{روند پیاده‌سازی}

پیاده‌سازی مدار طی مراحل زیر انجام شد:

\begin{enumerate}
\item ابتدا دو ثبات \lr{A} و \lr{B} با استفاده از دو تراشه‌ی \lr{74175} ساخته شدند. ورودی ساعت \lr{(clock)} هر دو ثبات به یک کلید فشاری (\lr{push-button}) مشترک متصل شد تا با هر بار فشردن، یک لبه‌ی ساعت به هر دو ثبات اعمال شود. همچنین ورودی \lr{Reset} هر دو ثبات به یک کلید فشاری دیگر برای بازگرداندن مدار به حالت اولیه وصل گردید.
\item در ادامه دو مالتی‌پلکسر با استفاده از تراشه‌ی \lr{74157} در مسیر ورودی هر ثبات قرار داده شدند تا امکان انتخاب بین داده‌ی خارجی و مقدار بازخوردی فراهم شود.
\item سپس تراشه‌ی \lr{74181} به‌عنوان \lr{ALU} به خروجی ثبات‌های \lr{A} و \lr{B} متصل شد و خروجی \lr{F} آن هم به‌عنوان خروجی نهایی مدار (\lr{F}) و هم به‌صورت بازخورد به مالتی‌پلکسر ورودی ثبات \lr{A} وصل گردید.
\item در نهایت، مدار کنترل‌کننده بر اساس جدول درستی \ref{tab:controller} با گیت‌های منطقی پایه پیاده‌سازی و به ورودی‌های \lr{SA}، \lr{SB}، \lr{S3-S0} و \lr{M} متصل شد. سه کلید (یا سوییچ) نیز برای اعمال دستی مقادیر \lr{M2-M0} در نظر گرفته شد.
\item پس از اتمام سیم‌کشی، برای هر یک از هشت حالت جدول \ref{tab:ops}، مقادیر \lr{M2-M0} به‌صورت دستی تنظیم و پس از اعمال چند مقدار نمونه در ورودی‌های \lr{D0-D3} و فشردن کلید ساعت، صحت مقدار ذخیره‌شده در ثبات‌ها و خروجی \lr{ALU} با استفاده از پروب‌های ولتاژ (\lr{voltage probe}) در \lr{Proteus} بررسی شد.
\end{enumerate}

در حین پیاده‌سازی، مهم‌ترین نکته‌ی موردتوجه، هم‌زمانی صحیح سیگنال‌های انتخاب (\lr{SA}، \lr{SB} و \lr{S3-S0}) با لبه‌ی ساعت بود؛ زیرا این سیگنال‌ها باید پیش از رسیدن لبه‌ی ساعت به ثبات‌ها، در مالتی‌پلکسرها اعمال شده باشند تا مقدار درستی وارد ثبات شود.

\section{دیاگرام یا شماتیک طرح اولیه}

طرح اولیه‌ی پیشنهادی، همان ساختار بلوک‌دیاگرام ارائه‌شده در راهنمای آزمایش بود: دو مالتی‌پلکسر ۴ بیتی در ورودی ثبات‌های \lr{A} و \lr{B}، خود دو ثبات \lr{A} و \lr{B}، واحد \lr{ALU} و بلوک کنترل‌کننده، بدون تعیین جزئیات داخلی گیت‌های کنترل‌کننده. در این مرحله صرفاً اتصالات کلی بین بلوک‌ها (مسیر داده‌ی خارجی \lr{D3-D0}، مسیر بازخورد از خروجی \lr{ALU} به مالتی‌پلکسر \lr{A}، و اتصال ثبات‌ها به ورودی‌های \lr{ALU}) مشخص شده بود و طراحی داخلی مدار کنترل‌کننده (یعنی گیت‌هایی که \lr{SA}، \lr{SB} و \lr{S3-S0} را تولید می‌کنند) در مرحله‌ی بعد، بر اساس جدول \ref{tab:controller}، تکمیل شد. این ساختار کلی همان چیزی است که در شکل \ref{fig:main} (شماتیک نهایی پیاده‌سازی‌شده) نیز قابل مشاهده است، با این تفاوت که در آن شکل، مدار کنترل‌کننده به‌طور کامل سیم‌کشی و پیاده‌سازی شده است.

\section{دیاگرام طرح نهایی}

پس از پیاده‌سازی اولیه، در حین تست مدار، ساختار کلی طرح (تعداد و نوع مالتی‌پلکسرها، ثبات‌ها و \lr{ALU}) بدون تغییر باقی ماند و صرفاً جزئیات داخلی مدار کنترل‌کننده (یعنی گیت‌های تولیدکننده‌ی \lr{SA}، \lr{SB} و \lr{S3-S0} طبق جدول \ref{tab:controller}) تکمیل و سیم‌کشی نهایی شد. به‌عبارت دیگر، طرح نهایی از نظر بلوک‌دیاگرام کلی با طرح اولیه یکسان است و تغییری در معماری کلی مدار اعمال نشده است؛ شماتیک نهایی همان شکل \ref{fig:main} است.

\section{بررسی \lr{Test Case} ها}

پس از تکمیل پیاده‌سازی، مدار برای تمامی هشت حالت ممکن ورودی \lr{M2-M0} (طبق جدول \ref{tab:ops}) آزمایش شد. برای هر حالت، مقادیر مناسبی در ورودی‌های \lr{D0-D3} اعمال و با فشردن کلید ساعت، نتیجه در خروجی ثبات‌ها و پروب‌های \lr{ALU} بررسی گردید. در حالت‌های \lr{101}، \lr{110} و \lr{111} که عملیات آن‌ها به مقدار اولیه‌ی هر دو ثبات \lr{A} و \lr{B} وابسته است، دو نمونه‌ی متفاوت آزمایش شد تا صحت عملکرد برای چند ورودی مختلف تأیید شود.

\subsection*{حالت \lr{$M2M1M0 = 000$}: \lr{$A \leftarrow D_3D_2D_1D_0$}}
در این حالت مقدار اعمال‌شده در خطوط \lr{D0-D3} پس از فشردن کلید ساعت، مستقیماً در ثبات \lr{A} ذخیره شد؛ همان‌طور که در شکل \ref{fig:t000} قابل مشاهده است.
\begin{figure}[H]
\centering
\includegraphics[width=0.8\textwidth]{5-1-test-000.PNG}
\caption{نتیجه‌ی تست حالت \lr{000}}
\label{fig:t000}
\end{figure}

\subsection*{حالت \lr{$M2M1M0 = 001$}: \lr{$B \leftarrow D_3D_2D_1D_0$}}
در این حالت مقدار خطوط \lr{D0-D3} در ثبات \lr{B} بارگذاری شد و ثبات \lr{A} بدون تغییر باقی ماند (شکل \ref{fig:t001}).
\begin{figure}[H]
\centering
\includegraphics[width=0.8\textwidth]{5-1-test-001.PNG}
\caption{نتیجه‌ی تست حالت \lr{001}}
\label{fig:t001}
\end{figure}

\subsection*{حالت \lr{$M2M1M0 = 010$}: \lr{$A \leftarrow A$}}
در این حالت مقدار ثبات \lr{A} بدون تغییر باقی ماند و پس از فشردن کلید ساعت نیز همان مقدار قبلی در \lr{A} حفظ شد (شکل \ref{fig:t010}).
\begin{figure}[H]
\centering
\includegraphics[width=0.8\textwidth]{5-1-test-010.PNG}
\caption{نتیجه‌ی تست حالت \lr{010}}
\label{fig:t010}
\end{figure}

\subsection*{حالت \lr{$M2M1M0 = 011$}: \lr{$A \leftarrow B$}}
مقدار ذخیره‌شده در ثبات \lr{B} در این حالت به ثبات \lr{A} منتقل شد (شکل \ref{fig:t011}).
\begin{figure}[H]
\centering
\includegraphics[width=0.8\textwidth]{5-1-test-011.PNG}
\caption{نتیجه‌ی تست حالت \lr{011}}
\label{fig:t011}
\end{figure}

\subsection*{حالت \lr{$M2M1M0 = 100$}: \lr{$clear(A)$}}
در این حالت با فشردن کلید ساعت، ثبات \lr{A} صفر شد (شکل \ref{fig:t100}).
\begin{figure}[H]
\centering
\includegraphics[width=0.8\textwidth]{5-1-test-100.PNG}
\caption{نتیجه‌ی تست حالت \lr{100}}
\label{fig:t100}
\end{figure}

\subsection*{حالت \lr{$M2M1M0 = 101$}: \lr{$A \leftarrow not(A)$}}
مقدار متمم ثبات \lr{A} در آن ذخیره شد. این حالت برای دو مقدار اولیه‌ی متفاوت آزمایش شد (شکل \ref{fig:t101}).
\begin{figure}[H]
\centering
\begin{subfigure}[t]{0.48\textwidth}
\centering
\includegraphics[width=\textwidth]{5-1-test-101-1.PNG}
\caption{نمونه‌ی اول}
\label{fig:t101-1}
\end{subfigure}
\hfill
\begin{subfigure}[t]{0.48\textwidth}
\centering
\includegraphics[width=\textwidth]{5-1-test-101-2.PNG}
\caption{نمونه‌ی دوم}
\label{fig:t101-2}
\end{subfigure}
\caption{نتیجه‌ی تست حالت \lr{101}}
\label{fig:t101}
\end{figure}

\subsection*{حالت \lr{$M2M1M0 = 110$}: \lr{$A \leftarrow and(A,B)$}}
حاصل \lr{AND} منطقی بین \lr{A} و \lr{B} در ثبات \lr{A} ذخیره شد. این حالت نیز برای دو نمونه‌ی مختلف بررسی شد (شکل \ref{fig:t110}).
\begin{figure}[H]
\centering
\begin{subfigure}[t]{0.48\textwidth}
\centering
\includegraphics[width=\textwidth]{5-1-test-110-1.PNG}
\caption{نمونه‌ی اول}
\label{fig:t110-1}
\end{subfigure}
\hfill
\begin{subfigure}[t]{0.48\textwidth}
\centering
\includegraphics[width=\textwidth]{5-1-test-110-2.PNG}
\caption{نمونه‌ی دوم}
\label{fig:t110-2}
\end{subfigure}
\caption{نتیجه‌ی تست حالت \lr{110}}
\label{fig:t110}
\end{figure}

\subsection*{حالت \lr{$M2M1M0 = 111$}: \lr{$A \leftarrow add(A,B)$}}
حاصل جمع محاسباتی \lr{A} و \lr{B} در ثبات \lr{A} ذخیره شد. این حالت نیز برای دو نمونه‌ی ورودی مختلف، ازجمله بررسی تولید یا عدم تولید رقم نقلی، آزمایش شد (شکل \ref{fig:t111}).
\begin{figure}[H]
\centering
\begin{subfigure}[t]{0.48\textwidth}
\centering
\includegraphics[width=\textwidth]{5-1-test-111-1.PNG}
\caption{نمونه‌ی اول}
\label{fig:t111-1}
\end{subfigure}
\hfill
\begin{subfigure}[t]{0.48\textwidth}
\centering
\includegraphics[width=\textwidth]{5-1-test-111-2.PNG}
\caption{نمونه‌ی دوم}
\label{fig:t111-2}
\end{subfigure}
\caption{نتیجه‌ی تست حالت \lr{111}}
\label{fig:t111}
\end{figure}

نتایج به‌دست‌آمده از تمامی هشت حالت فوق، کاملاً منطبق بر جدول عملیات \ref{tab:ops} بود و هیچ‌گونه ناسازگاری‌ای بین رفتار موردانتظار و رفتار مشاهده‌شده در شبیه‌سازی \lr{Proteus} مشاهده نشد.

\section{نتیجه‌گیری}

در این آزمایش، با ساخت و شبیه‌سازی مداری شامل تراشه‌ی \lr{74181} (\lr{ALU})، دو تراشه‌ی \lr{74175} (ثبات) و دو تراشه‌ی \lr{74157} (مالتی‌پلکسر)، نحوه‌ی کنترل عملیات محاسباتی و منطقی روی دو ثبات به‌صورت عملی بررسی شد. مهم‌ترین نکته‌ی آموخته‌شده، اهمیت طراحی صحیح مدار کنترل‌کننده برای تبدیل یک کد دستور کوتاه (سه‌بیتی) به سیگنال‌های انتخاب مناسب \lr{ALU} و مالتی‌پلکسرهاست؛ این ساختار در واقع نمونه‌ای ساده از معماری بخش کنترل و مسیر داده (\lr{control unit} و \lr{datapath}) در پردازنده‌های واقعی است.

آزمایش تمامی هشت حالت جدول عملیات نشان داد که مدار طراحی‌شده به‌درستی و مطابق انتظار عمل می‌کند و خروجی‌های به‌دست‌آمده در تمام موارد با مقادیر پیش‌بینی‌شده مطابقت دارد. این آزمایش درک بهتری از کاربرد عملی تراشه‌ی \lr{74181} به‌عنوان یک واحد محاسبات و منطق کلاسیک، و نیز نحوه‌ی طراحی مدارهای کنترل‌کننده‌ی ترکیبی برای سیستم‌های دیجیتال کوچک فراهم کرد.

\newpage
\part{بخش دوم: طراحی \lr{ALU} چهاربیتی با تراشه‌های \lr{TTL} پایه}

\section{مقدمه و تعریف مسئله}

هدف از آزمایش پنجم، آشنایی عملی با نحوه‌ی طراحی و پیاده‌سازی یک \textbf{واحد محاسبات و منطق} (\en{Arithmetic Logic Unit} یا به‌اختصار \en{ALU}) است. \en{ALU} یکی از اصلی‌ترین اجزای هر واحد پردازش دیجیتال (مانند پردازنده‌های مرکزی) محسوب می‌شود که وظیفه‌ی انجام عملیات حسابی (جمع، تفریق، افزایش، کاهش و مانند آن) و عملیات منطقی (\en{AND}،‌ \en{OR}، \en{XOR} و متمم) و همچنین عملیات جابه‌جایی بیت‌ها (شیفت) را بر عهده دارد. تمامی مراحل طراحی، شبیه‌سازی و آزمایش این مدار با استفاده از نرم‌افزار \en{Proteus} انجام شده است.

مسئله‌ای که در این آزمایش باید حل شود به این صورت تعریف می‌شود: یک \en{ALU} چهار‌بیتی طراحی و پیاده‌سازی شود که دارای ورودی‌های زیر باشد:

\begin{itemize}[itemsep=2pt]
    \item دو عملوند چهاربیتی \en{A (A0-A3)} و \en{B (B0-B3)}
    \item یک ورودی کری \en{Cin}
    \item چهار خط انتخاب عملیات \en{Select (S0-S3)}
\end{itemize}

و دارای خروجی‌های زیر باشد:

\begin{itemize}[itemsep=2pt]
    \item یک خروجی چهاربیتی \en{F (F0-F3)}
    \item یک خروجی کری خروجی \en{Cout}
\end{itemize}

با توجه به مقدار خطوط انتخاب \en{S3 S2 S1 S0} و ورودی \en{Cin}، مدار باید یکی از چهارده عملیات مشخص‌شده در جدول عملکرد (جدول \ref{tab:optable}) را روی عملوندهای ورودی اجرا کرده و نتیجه را در خروجی \en{F} (و در حالت عملیات حسابی، کری خروجی را در \en{Cout}) قرار دهد. در پایان آزمایش نیز باید صحت عملکرد مدار برای تمامی حالت‌های ممکنِ خط انتخاب، توسط ابزار \en{Logic Probe} در محیط \en{Proteus} بررسی و مستند شود.

\section{بررسی و انتخاب روش طراحی}

برای ساخت مدار داخلی یک \en{ALU} (که یک واحد محاسبات و منطق است) به‌جای طراحی مدار کاملاً سفارشی از دروازه‌های منطقی پایه، از رویکرد \textbf{استفاده از تراشه‌های \en{TTL} آماده} بهره گرفته شده است؛ چراکه این روش هم پیاده‌سازی را ساده‌تر می‌کند و هم با تجهیزات آزمایشگاهی و کتابخانه‌ی \en{Proteus} سازگاری کامل دارد. اجزای اصلی به‌کاررفته عبارت‌اند از:

\begin{itemize}[itemsep=2pt]
    \item یک تراشه \en{74283} (جمع‌کننده‌ی کامل چهاربیتی) به‌عنوان هسته‌ی محاسباتیِ جمع/تفریق؛
    \item شش تراشه \en{74153} (مالتی‌پلکسر دوگانه‌ی چهار‌به‌یک) برای انتخاب عملوندها و ترکیب نتایج بخش‌های مختلف؛
    \item دروازه‌های \en{7408} (\en{AND})، \en{7432} (\en{OR})، \en{74LS86} (\en{XOR}) و \en{7404} (\en{NOT}) به تعداد کافی برای پیاده‌سازی واحد منطقی و تولید سیگنال‌های کمکی (مانند \lr{$\overline{B}$} و \lr{$\overline{A}$}).
\end{itemize}

دلیل انتخاب این روش آن است که تراشه‌ی \en{74283} به‌طور مستقیم عملیات جمع چهاربیتی همراه با کری ورودی/خروجی را انجام می‌دهد و با تغذیه‌ی مناسب ورودی دوم آن (به‌جای \en{B} ساده) می‌توان تمامی هشت عملیات حسابی جدول را -- بدون نیاز به مدار جمع/تفریق‌کننده‌ی جداگانه -- تنها با انتخاب مقدار مناسب برای این ورودی (توسط مالتی‌پلکسر) پیاده‌سازی کرد. همچنین استفاده از مالتی‌پلکسرهای \en{74153} در سطوح بعدی، امکان می‌دهد که سه بخشِ مستقل مدار (حسابی، منطقی و شیفت) به‌صورت جداگانه طراحی شوند و در نهایت با یک لایه‌ی انتخابِ نهایی، بر اساس خطوط \en{S3} و \en{S2}، خروجی مناسب انتخاب گردد. این تفکیک، طراحی، اشکال‌زدایی و سیم‌کشی مدار را به‌مراتب ساده‌تر از یک پیاده‌سازی یکپارچه می‌کند.

به‌طور کلی مدار به سه بخش اصلی تقسیم می‌شود:

\begin{enumerate}[itemsep=2pt]
    \item واحد جمع و تفریق‌کننده (عملیات حسابی)
    \item واحد منطق پایه (\en{AND}، \en{OR}، \en{XOR}، متمم)
    \item واحد شیفت به راست و چپ
\end{enumerate}

که خروجی هر یک از این سه بخش، به‌همراه سیگنال‌های انتخاب \en{S3} و \en{S2}، به مالتی‌پلکسرهای نهایی داده می‌شود تا در نهایت یکی از آن‌ها به‌عنوان خروجی \en{F} انتخاب شود. جزئیات کامل هر بخش در قسمت «روند پیاده‌سازی» آمده است.

\section{طرح پیشنهادی اولیه}

پیش از ورود به جزئیات مدار و سیم‌کشی نهایی، ساختار کلی سیستم به‌صورت مفهومی و بر پایه‌ی تحلیل نظری تعریف شد. طبق این طرح اولیه، \en{ALU} چهاربیتی باید دارای دو ورودی عملوند چهاربیتی \en{A} و \en{B}، یک ورودی کری \en{Cin} و یک ورودی انتخاب چهاربیتی \en{Select} باشد و در خروجی، نتیجه‌ی چهاربیتی \en{F} به‌همراه کری خروجی \en{Cout} تولید کند. از آنجا که در این مرحله صرفاً ساختار کلی و تقسیم‌بندی وظایف مدار مشخص می‌شد -- و نه جزئیات سیم‌کشی یا چیدمان قطعات در محیط \en{Proteus} -- هیچ شماتیک یا بلوک‌دیاگرام تصویری از این مرحله تهیه نشده است؛ آنچه در ادامه می‌آید، شرح متنی همین طرح مفهومی اولیه است.

بر اساس این تحلیل اولیه، مقرر شد مدار داخلی \en{ALU} به سه زیرواحدِ مستقل تقسیم شود:

\begin{enumerate}[itemsep=2pt]
    \item زیرواحد حسابی، مسئول انجام هشت عملیات جمع، تفریق، افزایش و کاهش؛
    \item زیرواحد منطقی، مسئول انجام چهار عملیات \en{AND}، \en{OR}، \en{XOR} و متمم؛
    \item زیرواحد شیفت، مسئول جابه‌جایی بیت‌های \en{A} به راست یا چپ.
\end{enumerate}

و در نهایت، خروجی هر یک از این سه زیرواحد، توسط یک لایه‌ی مالتی‌پلکسرِ انتخاب نهایی -- بر اساس مقدار خطوط \en{S3} و \en{S2} -- به‌عنوان خروجی \en{F} مدار انتخاب گردد. این تقسیم‌بندی مفهومی، مبنای طراحی و سیم‌کشی مدار در بخش «روند پیاده‌سازی» قرار گرفت.

جدول عملکرد کامل که مبنای طراحی این آزمایش قرار گرفته، در جدول \ref{tab:optable} آمده است. همان‌طور که دیده می‌شود، ترکیب خطوط انتخاب \en{S3 S2} تعیین می‌کند که مدار در کدام‌یک از سه حالت (حسابی، منطقی یا شیفت) قرار گیرد و در ادامه، \en{S1 S0} (و در حالت حسابی، \en{Cin} نیز) عملیات دقیق درون آن حالت را مشخص می‌کنند.

\begin{table}[H]
\centering
\caption{جدول عملکرد \en{ALU} چهاربیتی (مبنای طراحی)}
\label{tab:optable}
\renewcommand{\arraystretch}{1.4}
\rowcolors{2}{rowlightB}{white}
\begin{tabular}{>{\centering\arraybackslash}p{0.7cm} >{\centering\arraybackslash}p{0.7cm} >{\centering\arraybackslash}p{0.7cm} >{\centering\arraybackslash}p{0.7cm} >{\centering\arraybackslash}p{0.9cm} | >{\centering\arraybackslash}p{3.4cm} >{\centering\arraybackslash}p{3cm}}
\toprule
\rowcolor{headerbgB}
\textcolor{white}{\lr{$S_3$}} & \textcolor{white}{\lr{$S_2$}} & \textcolor{white}{\lr{$S_1$}} & \textcolor{white}{\lr{$S_0$}} & \textcolor{white}{\lr{$C_{in}$}} & \textcolor{white}{عملیات} & \textcolor{white}{توضیح} \\
\midrule
\en{0} & \en{0} & \en{0} & \en{0} & \en{0} & \en{$F = A$}                     & انتقال \en{A} \\
\en{0} & \en{0} & \en{0} & \en{0} & \en{1} & \en{$F = A + 1$}                 & افزایش (اینکریمنت) \en{A} \\
\en{0} & \en{0} & \en{0} & \en{1} & \en{0} & \en{$F = A + B$}                 & جمع \\
\en{0} & \en{0} & \en{0} & \en{1} & \en{1} & \en{$F = A + B + 1$}             & جمع با کری \\
\en{0} & \en{0} & \en{1} & \en{0} & \en{0} & \en{$F = A + \overline{B}$}      & تفریق با قرض \\
\en{0} & \en{0} & \en{1} & \en{0} & \en{1} & \en{$F = A + \overline{B} + 1$}  & تفریق \\
\en{0} & \en{0} & \en{1} & \en{1} & \en{0} & \en{$F = A - 1$}                 & کاهش (دیکرمنت) \en{A} \\
\en{0} & \en{0} & \en{1} & \en{1} & \en{1} & \en{$F = A$}                     & انتقال \en{A} \\
\en{0} & \en{1} & \en{0} & \en{0} & \en{$\times$} & \en{$F = A \wedge B$}     & \en{AND} \\
\en{0} & \en{1} & \en{0} & \en{1} & \en{$\times$} & \en{$F = A \vee B$}       & \en{OR} \\
\en{0} & \en{1} & \en{1} & \en{0} & \en{$\times$} & \en{$F = A \oplus B$}     & \en{XOR} \\
\en{0} & \en{1} & \en{1} & \en{1} & \en{$\times$} & \en{$F = \overline{A}$}   & متمم \en{A} \\
\en{1} & \en{0} & \en{$\times$} & \en{$\times$} & \en{$\times$} & \en{$F = \mathrm{shr}\ A$} & شیفت به راستِ \en{A} در \en{F} \\
\en{1} & \en{1} & \en{$\times$} & \en{$\times$} & \en{$\times$} & \en{$F = \mathrm{shl}\ A$} & شیفت به چپِ \en{A} در \en{F} \\
\bottomrule
\end{tabular}
\end{table}

\section{روند پیاده‌سازی}

پیاده‌سازی مدار در سه گام مطابق با تفکیک طرح اولیه انجام شد.

\subsection{بخش اول: واحد جمع و تفریق‌کننده}

در این بخش رابطه‌ی کلی زیر پیاده‌سازی می‌شود:
\begin{center}
\en{$F = A + X + C_{in}$}
\end{center}
که در آن \en{X} می‌تواند \en{B}، متمم \en{B} (یعنی \en{$\overline{B}$})، مقدار صفر یا مقدار یک (تمام بیت‌ها) باشد. برای انتخاب این چهار حالت از دو تراشه‌ی مالتی‌پلکسر \en{74153} استفاده شده است (از آنجا که ورودی‌ها و خروجی‌ها چهاربیتی هستند و هر تراشه‌ی \en{74153} از دو مالتی‌پلکسر چهار‌به‌یک تشکیل شده، به دو تراشه یعنی چهار واحد مالتی‌پلکسر نیاز است تا هر بیت از \en{X} به‌طور مستقل انتخاب شود).

پایه‌های ورودی هر مالتی‌پلکسر به‌ترتیب به تغذیه (\en{VCC})، \en{B}، متمم \en{B} و زمین (\en{GND}) متصل شده‌اند. سپس بیت‌های \en{A0} تا \en{A3} به ورودی‌های \en{A0} تا \en{A3} تراشه‌ی \en{74283} و بیت‌های \en{B0} تا \en{B3} این تراشه به خروجی‌های چهار مالتی‌پلکسر (یعنی مقدار انتخاب‌شده‌ی \en{X}) متصل می‌شوند. پایه‌های انتخاب مالتی‌پلکسرها نیز به \en{S0} و \en{S1} وصل شده‌اند تا مطابق جدول عملکرد، یکی از چهار حالت \en{X} انتخاب شود.

در خروجی تراشه‌ی جمع‌کننده، کری خروجی (\en{Cout}) مستقیماً به \en{Logic Probe} متصل می‌شود، اما خروجی‌های اصلی جمع (\en{F0} تا \en{F3}) باید به ورودی اول دو مالتی‌پلکسر نهایی (که در بخش بعد معرفی می‌شوند) داده شوند تا در انتخاب نهایی خروجی \en{F} شرکت کنند.

\subsection{بخش دوم: واحد منطقی}

در این بخش ابتدا چهار عملیات پایه‌ی \en{$A \wedge B$}، \en{$A \vee B$}، \en{$A \oplus B$} و \en{$\overline{A}$} با استفاده از گیت‌های \en{7408}، \en{7432}، \en{74LS86} و \en{7404} ساخته می‌شوند. سپس با استفاده از دو تراشه‌ی مالتی‌پلکسر دیگر، یکی از این چهار نتیجه انتخاب می‌شود: ورودی اول هر مالتی‌پلکسر به خروجی گیت \en{AND}، ورودی دوم به \en{OR}، ورودی سوم به \en{XOR} و ورودی چهارم به گیت \en{NOT} (متمم \en{A}) متصل است. پایه‌های انتخاب این مالتی‌پلکسرها نیز مانند بخش قبل به \en{S0} و \en{S1} وصل شده‌اند. خروجی این دو مالتی‌پلکسر به ورودی دوم مالتی‌پلکسرهای نهایی داده می‌شود.

\subsection{بخش سوم: واحد شیفت}

برای این بخش نیازی به گیت اضافه نیست و شیفت به‌صورت مستقیم با سیم‌کشیِ مناسب ورودی‌های \en{A} پیاده‌سازی می‌شود؛ برای شیفت به راست، بیت \en{A0} به زمین وصل شده و سایر بیت‌ها به‌ترتیب یک واحد جابه‌جا می‌شوند و به همین ترتیب برعکس، برای شیفت به چپ. این سیم‌کشی مستقیماً به ورودی سوم و چهارم مالتی‌پلکسرهای نهایی متصل می‌شود.

پایه‌های انتخاب دو مالتی‌پلکسر نهایی به \en{S2} و \en{S3} وصل شده‌اند تا در نهایت یکی از چهار مسیر (حسابی، منطقی، شیفت راست، شیفت چپ) به‌عنوان خروجی \en{F} انتخاب و به \en{Logic Probe}‌های خروجی نمایش داده شود.

با توجه به این‌که پس از پیاده‌سازی، طرح اولیه به‌طور کامل پاسخ‌گوی نیاز آزمایش بود، \textbf{نسخه‌ی اولیه و نهایی این مدار هیچ تفاوتی با یکدیگر نداشتند} و نیازی به هیچ تغییری در ساختار یا سیم‌کشی مدار پیش نیامد.

\section{شماتیک و تصویر مدار نهایی}

لازم به ذکر است که در این آزمایش، طراحی مدار مستقیماً بر پایه‌ی طرح مفهومی اولیه (بخش «طرح پیشنهادی اولیه») و جدول عملکرد جدول \ref{tab:optable} انجام شد و در همان تلاش نخست، بدون نیاز به هیچ بازبینی یا اصلاح، به‌درستی پاسخ‌گوی تمامی حالت‌ها بود؛ به‌عبارت دیگر، این مدار اساساً \textbf{مرحله یا نسخه‌ی «اولیه»ی جداگانه‌ای (متفاوت از طرح نهایی) نداشته است}. به همین دلیل هیچ تصویر یا شماتیکی از یک نسخه‌ی «اولیه» در دسترس نیست؛ آنچه در ادامه ارائه می‌شود -- هم شماتیک و هم تصویر واقعی مدار -- هر دو متعلق به همان \textbf{یک نسخه‌ی طراحی‌شده و نهایی} مدار، در محیط \en{Proteus}، هستند.

شماتیک کامل مدار نهایی، شامل تراشه‌ی جمع‌کننده‌ی \en{74283}\ (با نام \en{U4})، شش مالتی‌پلکسر \en{74153}\ (\en{U1}، \en{U2}، \en{U9}، \en{U10}، \en{U11} و \en{U12})، دروازه‌های \en{AND}\ (\en{U5})، \en{OR}\ (\en{U6})، \en{XOR}\ (\en{U13} و \en{U7}) و معکوس‌کننده‌های \en{NOT}\ (\en{U3} و \en{U8})، در شکل \ref{fig:schematic} نشان داده شده است.

\begin{figure}[H]
    \centering
    \includegraphics[width=\textwidth]{main_schematic_cropped.pdf}
    \caption{شماتیک کامل مدار نهاییِ \en{ALU} چهاربیتی در محیط \en{Proteus}}
    \label{fig:schematic}
\end{figure}

همان‌طور که در شکل مشخص است، ورودی‌های \en{A0-A3}، \en{B0-B3}، \en{S0-S3} و \en{Cin} در سمت چپ مدار قرار دارند و خروجی‌های \en{F0-F3} و \en{Cout} در سمت راست مدار نمایش داده می‌شوند. ساختار سه‌بخشی مدار (واحد حسابی در بالا، واحد منطقی در وسط و واحد شیفت به‌صورت سیم‌کشی مستقیم به مالتی‌پلکسرهای نهایی) در این شماتیک کاملاً قابل‌مشاهده است.

\section{بررسی \en{Test Case}‌ها}

برای بررسی صحت عملکرد مدار، تمامی چهارده حالت جدول عملکرد (جدول \ref{tab:optable}) به‌صورت جداگانه در \en{Proteus} تحریک شدند و مقادیر \en{Logic Probe}‌های متصل به \en{F0-F3} و \en{Cout} برای هر حالت ثبت شد. نام‌گذاری فایل‌های تصویر هر \en{Test Case} بر اساس ترتیب \en{$S_3 S_2 S_1 S_0 C_{in}$} انجام شده است (علامت \en{x} به‌معنای بی‌اثر بودن آن بیت در عملیات موردنظر است). جدول \ref{tab:testcases} نگاشت میان نام هر فایل تصویر و عملیات متناظر با آن را نشان می‌دهد.

\begin{table}[H]
\centering
\caption{نگاشت \en{Test Case}‌ها به عملیات جدول عملکرد}
\label{tab:testcases}
\renewcommand{\arraystretch}{1.4}
\rowcolors{2}{rowlightB}{white}
\begin{tabular}{>{\centering\arraybackslash}p{3cm} >{\centering\arraybackslash}p{5.5cm}}
\toprule
\rowcolor{headerbgB}
\textcolor{white}{\lr{$S_3S_2S_1S_0C_{in}$}} & \textcolor{white}{عملیات} \\
\midrule
\en{0 0 0 0 0} & \en{$F=A$} (انتقال \en{A}) \\
\en{0 0 0 0 1} & \en{$F=A+1$} (افزایش) \\
\en{0 0 0 1 0} & \en{$F=A+B$} (جمع) \\
\en{0 0 0 1 1} & \en{$F=A+B+1$} (جمع با کری) \\
\en{0 0 1 0 0} & \en{$F=A+\overline{B}$} (تفریق با قرض) \\
\en{0 0 1 0 1} & \en{$F=A+\overline{B}+1$} (تفریق) \\
\en{0 0 1 1 0} & \en{$F=A-1$} (کاهش) \\
\en{0 0 1 1 1} & \en{$F=A$} (انتقال \en{A}) \\
\en{0 1 0 0 x} & \en{$F=A\wedge B$} (\en{AND}) \\
\en{0 1 0 1 x} & \en{$F=A\vee B$} (\en{OR}) \\
\en{0 1 1 0 x} & \en{$F=A\oplus B$} (\en{XOR}) \\
\en{0 1 1 1 x} & \en{$F=\overline{A}$} (متمم \en{A}) \\
\en{1 0 x x x} & \en{$F=\mathrm{shr}\ A$} (شیفت راست) \\
\en{1 1 x x x} & \en{$F=\mathrm{shl}\ A$} (شیفت چپ) \\
\bottomrule
\end{tabular}
\end{table}

پیش از بررسی تک‌تک حالت‌ها، نمای کلی مدار نهاییِ پیاده‌سازی‌شده به‌همراه محل قرارگیری کلیدهای ورودی و \en{Logic Probe}‌های خروجی پیش‌تر در شکل \ref{fig:schematic} (بخش «شماتیک و تصویر مدار نهایی») نشان داده شد.

\subsection{عملیات حسابی}

هشت حالت اول جدول عملکرد مربوط به عملیات حسابی هستند که با \en{$S_3S_2=00$} و ترکیب‌های مختلف \en{$S_1S_0C_{in}$} فعال می‌شوند. نتایج شبیه‌سازی برای این هشت حالت در شکل‌های \ref{fig:arith1} و \ref{fig:arith2} نشان داده شده‌اند.

\begin{figure}[H]
    \centering
    \begin{subfigure}[t]{0.48\textwidth}
        \centering
        \includegraphics[width=\textwidth]{00000.png}
        \caption{\en{$F=A$}}
    \end{subfigure}
    \hfill
    \begin{subfigure}[t]{0.48\textwidth}
        \centering
        \includegraphics[width=\textwidth]{00001.png}
        \caption{\en{$F=A+1$}}
    \end{subfigure}

    \vspace{0.4cm}
    \begin{subfigure}[t]{0.48\textwidth}
        \centering
        \includegraphics[width=\textwidth]{00010.png}
        \caption{\en{$F=A+B$}}
    \end{subfigure}
    \hfill
    \begin{subfigure}[t]{0.48\textwidth}
        \centering
        \includegraphics[width=\textwidth]{00011.png}
        \caption{\en{$F=A+B+1$}}
    \end{subfigure}
    \caption{نتیجه‌ی شبیه‌سازی چهار عملیات حسابی نخست}
    \label{fig:arith1}
\end{figure}

\begin{figure}[H]
    \centering
    \begin{subfigure}[t]{0.48\textwidth}
        \centering
        \includegraphics[width=\textwidth]{00100.png}
        \caption{\en{$F=A+\overline{B}$}}
    \end{subfigure}
    \hfill
    \begin{subfigure}[t]{0.48\textwidth}
        \centering
        \includegraphics[width=\textwidth]{00101.png}
        \caption{\en{$F=A+\overline{B}+1$}}
    \end{subfigure}

    \vspace{0.4cm}
    \begin{subfigure}[t]{0.48\textwidth}
        \centering
        \includegraphics[width=\textwidth]{00110.png}
        \caption{\en{$F=A-1$}}
    \end{subfigure}
    \hfill
    \begin{subfigure}[t]{0.48\textwidth}
        \centering
        \includegraphics[width=\textwidth]{00111.png}
        \caption{\en{$F=A$}}
    \end{subfigure}
    \caption{نتیجه‌ی شبیه‌سازی چهار عملیات حسابی دوم}
    \label{fig:arith2}
\end{figure}

در تمامی هشت حالت فوق، مقدار خوانده‌شده از \en{Logic Probe}‌های \en{F0-F3} با مقدار محاسبه‌شده‌ی دستی از روی جدول عملکرد یکسان بود و کری خروجی (\en{Cout}) نیز به‌درستی وضعیت سرریز جمع/تفریق را نشان می‌داد.

\subsection{عملیات منطقی}

چهار عملیات منطقی با \en{$S_3S_2=01$} و مقادیر مختلف \en{$S_1S_0$} فعال می‌شوند (مقدار \en{Cin} در این حالت بی‌اثر است). نتایج در شکل \ref{fig:logic} نشان داده شده‌اند.

\begin{figure}[H]
    \centering
    \begin{subfigure}[t]{0.48\textwidth}
        \centering
        \includegraphics[width=\textwidth]{0100x.png}
        \caption{\en{$F=A\wedge B$}\ (\en{AND})}
    \end{subfigure}
    \hfill
    \begin{subfigure}[t]{0.48\textwidth}
        \centering
        \includegraphics[width=\textwidth]{0101x.png}
        \caption{\en{$F=A\vee B$}\ (\en{OR})}
    \end{subfigure}

    \vspace{0.4cm}
    \begin{subfigure}[t]{0.48\textwidth}
        \centering
        \includegraphics[width=\textwidth]{0110x.png}
        \caption{\en{$F=A\oplus B$}\ (\en{XOR})}
    \end{subfigure}
    \hfill
    \begin{subfigure}[t]{0.48\textwidth}
        \centering
        \includegraphics[width=\textwidth]{0111x.png}
        \caption{\en{$F=\overline{A}$}\ (متمم)}
    \end{subfigure}
    \caption{نتیجه‌ی شبیه‌سازی چهار عملیات منطقی}
    \label{fig:logic}
\end{figure}

در هر چهار حالت، خروجی \en{F} دقیقاً برابر با نتیجه‌ی منطقی موردانتظار روی عملوندهای \en{A} و \en{B} بود و همان‌طور که از جدول عملکرد انتظار می‌رفت، مقدار \en{Cout} در این حالت فاقد معنای مشخص (بی‌اثر) است.

\subsection{عملیات شیفت}

دو عملیات پایانی جدول، شیفت به راست و شیفت به چپ \en{A} هستند که با \en{$S_3S_2=10$} و \en{$S_3S_2=11$} به‌ترتیب فعال می‌شوند. نتایج در شکل \ref{fig:shift} آمده است.

\begin{figure}[H]
    \centering
    \begin{subfigure}[t]{0.48\textwidth}
        \centering
        \includegraphics[width=\textwidth]{10xxx.png}
        \caption{\en{$F=\mathrm{shr}\ A$}}
    \end{subfigure}
    \hfill
    \begin{subfigure}[t]{0.48\textwidth}
        \centering
        \includegraphics[width=\textwidth]{11xxx.png}
        \caption{\en{$F=\mathrm{shl}\ A$}}
    \end{subfigure}
    \caption{نتیجه‌ی شبیه‌سازی دو عملیات شیفت}
    \label{fig:shift}
\end{figure}

در هر دو حالت، بیت جدیدی که از سمت آزاد وارد \en{F} می‌شود برابر با صفر بود و سایر بیت‌های \en{A} دقیقاً یک واحد به سمت مربوطه جابه‌جا شدند که کاملاً منطبق با تعریف عملیات شیفت منطقی در جدول عملکرد است.

\section{نتیجه‌گیری}

در این آزمایش یک واحد محاسبات و منطق (\en{ALU}) چهاربیتی، مطابق با جدول عملکرد استاندارد (چهارده عملیات حسابی، منطقی و شیفت)، با استفاده از تراشه‌ی جمع‌کننده‌ی \en{74283}، شش تراشه‌ی مالتی‌پلکسر \en{74153} و تعداد کافی گیت‌های \en{AND}، \en{OR}، \en{XOR} و \en{NOT} طراحی و در نرم‌افزار \en{Proteus} پیاده‌سازی شد.

تفکیک مدار به سه زیرواحد مستقلِ حسابی، منطقی و شیفت -- که در نهایت توسط یک لایه‌ی مالتی‌پلکسر انتخاب می‌شدند -- روش طراحی را به‌طور قابل‌توجهی ساده‌تر و قابل‌فهم‌تر کرد و نیاز به هیچ مدار جمع/تفریق‌کننده یا واحد منطقی جداگانه‌ای نبود؛ تمام عملیات حسابی صرفاً با تغییر ورودی دوم تراشه‌ی جمع‌کننده (توسط مالتی‌پلکسرهای سطح اول) به‌دست آمدند.

با اعمال تمامی چهارده \en{Test Case} ممکن (مطابق جدول \ref{tab:testcases})، خروجی مدار در همه‌ی حالت‌ها دقیقاً با مقدار موردانتظار از روی جدول عملکرد مطابقت داشت و نیازی به اصلاح طرح اولیه پیش نیامد؛ به همین دلیل طرح اولیه هیچ تفاوتی با طرح نهایی نداشت و شماتیک و تصویر ارائه‌شده در این گزارش، هر دو متعلق به همان مدار نهاییِ پیاده‌سازی‌شده هستند. مهم‌ترین نکته‌ی آموخته‌شده در این آزمایش، توانایی ترکیب چند تراشه‌ی \en{MSI} (نظیر جمع‌کننده و مالتی‌پلکسر) به‌جای طراحی مدار از پایه با دروازه‌های منطقی است؛ رویکردی که در طراحی واقعی پردازنده‌ها نیز به‌طور گسترده استفاده می‌شود.

\end{document}
